% Exemplo: revelar o conteúdo aos poucos no beamer: \pause, especificadores <...>, \uncover e \only
% Compilar: latexmk -pdf sobreposicoes.tex         Testado com: beamer, TeX Live 2024 ou mais recente
% Cada passo de um frame vira uma página do PDF.
\documentclass[aspectratio=169]{beamer}
\usepackage[T1]{fontenc}
\usepackage[brazilian]{babel}
\usetheme{Madrid}

\begin{document}

%%BEGIN:beamerpause
\begin{frame}{Comando pause}
  Os tópicos aparecem um de cada vez.
  \begin{itemize}
    \pause \item Primeiro tópico
    \pause \item Segundo tópico
  \end{itemize}
\end{frame}
%%END:beamerpause

%%BEGIN:beamerespecificadores
\begin{frame}{Especificadores de sobreposição}
  \begin{itemize}
    \item<1>  Aparece só no primeiro passo
    \item<2-3> Aparece no segundo e no terceiro
    \item<3-> Aparece a partir do terceiro e continua
  \end{itemize}
  \begin{block}<4->{Bloco}
    O bloco só aparece no quarto passo.
  \end{block}
\end{frame}
%%END:beamerespecificadores

%%BEGIN:beamerincremental
\begin{frame}{Lista incremental}
  \begin{itemize}[<+->]          % equivale a <+-> em cada \item
    \item Primeiro tópico
    \item Segundo tópico
    \item Terceiro tópico
  \end{itemize}
  \uncover<3->{Visível (e com espaço reservado) a partir do terceiro passo.}
  \only<2>{Só no segundo passo, sem reservar espaço nos outros.}
\end{frame}
%%END:beamerincremental

\end{document}
